\chapter{JADWAL PENELITIAN}

Penelitian ini direncanakan untuk dilaksanakan selama delapan bulan, dari bulan Juni 2026 hingga Januari 2027. Jadwal penelitian disusun berdasarkan tahapan evaluasi empiris terkontrol: persiapan lingkungan, pengembangan skrip otomasi, kalibrasi baseline, pelaksanaan eksperimen, analisis statistik, dan penulisan laporan.

Tabel~\ref{tab:jadwal} menunjukkan rencana jadwal pelaksanaan penelitian berbasis bulan. Simbol $\bullet$ menunjukkan bahwa kegiatan dilaksanakan pada bulan tersebut.

\begin{table}[H]
  \centering
  \caption{Jadwal Penelitian (Gantt Chart)}
  \label{tab:jadwal}
  \resizebox{\textwidth}{!}{%
  \begin{tabular}{|l|c|c|c|c|c|c|c|c|}
    \hline
    \multirow{2}{*}{\textbf{Kegiatan}} & \multicolumn{8}{c|}{\textbf{Bulan}} \\
    \cline{2-9}
     & \textbf{Jun} & \textbf{Jul} & \textbf{Agu} & \textbf{Sep} & \textbf{Okt} & \textbf{Nov} & \textbf{Des} & \textbf{Jan} \\
    \hline
    Studi Literatur & $\bullet$ & $\bullet$ & & & & & & \\
    \hline
    Perancangan Lingkungan & & $\bullet$ & $\bullet$ & & & & & \\
    \hline
    Implementasi ArgoCD + Baseline Manual & & & $\bullet$ & $\bullet$ & & & & \\
    \hline
    Pengembangan Skrip Otomasi & & $\bullet$ & $\bullet$ & $\bullet$ & & & & \\
    \hline
    Persiapan Baseline \& Kalibrasi & & & & $\bullet$ & $\bullet$ & & & \\
    \hline
    Eksperimen \& Pengumpulan Data (E1--E3) & & & & & $\bullet$ & $\bullet$ & & \\
    \hline
    Analisis Statistik \& Interpretasi & & & & & & $\bullet$ & $\bullet$ & \\
    \hline
    Penulisan Laporan & $\bullet$ & $\bullet$ & $\bullet$ & $\bullet$ & $\bullet$ & $\bullet$ & $\bullet$ & $\bullet$ \\
    \hline
    Seminar / Sidang & & & & & & & & $\bullet$ \\
    \hline
  \end{tabular}%
  }
\end{table}

\section{Penjelasan Kegiatan}

\begin{enumerate}
    \item \textbf{Studi Literatur (Juni -- Juli 2026)}
    
    Kegiatan ini mencakup pengumpulan dan analisis literatur mengenai: \textit{configuration drift} pada Kubernetes, evaluasi empiris GitOps, metode eksperimental pada riset sistem, dan arsitektur InvenioRDM sebagai beban uji. Sumber literatur meliputi jurnal ilmiah (IEEE, ACM, Springer), dokumentasi resmi, buku teks, dan artikel teknis. Tujuan dari fase ini adalah membangun landasan teoretis yang kuat.
    
    \item \textbf{Perancangan Lingkungan (Juli -- Agustus 2026)}
    
    Tahap ini meliputi perancangan arsitektur Lingkungan A (deployment manual/\texttt{kubectl}) dan Lingkungan B (deployment GitOps/ArgoCD) pada klaster Kubernetes yang sama. Perancangan meliputi: definisi variabel dan intervensi, pemilihan tujuh skenario drift, spesifikasi protokol pengukuran untuk tiga metrik (R1, R2, R3), dan rancangan eksperimen.
    
    \item \textbf{Implementasi ArgoCD + Baseline Manual (Agustus -- September 2026)}
    
    Tahap implementasi meliputi: konfigurasi ArgoCD pada klaster, pembuatan repositori GitOps dengan manifest deklaratif, konfigurasi baseline manual (\textit{workflow} \texttt{kubectl}), serta konfigurasi \textit{logging} dan \textit{monitoring} (Prometheus, Grafana, Kubernetes audit logs, ArgoCD notifications). Pada akhir fase ini, kedua lingkungan siap untuk eksperimen.
    
    \item \textbf{Pengembangan Skrip Otomasi (Juli -- September 2026)}
    
    Pengembangan empat skrip otomasi yang mendukung konsistensi prosedural eksperimen: (a)~\texttt{drift-inject.sh} untuk \textit{injection} skenario drift secara acak dengan \textit{logging} \textit{timestamp}; (b)~\texttt{metrics-collect.sh} untuk pengumpulan metrik secara berkala via ArgoCD API dan \texttt{kubectl}; (c)~\texttt{baseline-check.sh} untuk verifikasi kondisi klaster sebelum setiap \textit{run}; dan (d)~\texttt{recovery-runbook.sh} untuk memandu operator pada Lingkungan A melalui prosedur \textit{recovery} terstandarisasi. Skrip dikembangkan secara paralel dengan perancangan dan implementasi lingkungan agar siap digunakan saat eksperimen dimulai.
    
    \item \textbf{Persiapan Baseline \& Kalibrasi (September -- Oktober 2026)}
    
    Sebelum eksperimen dimulai, fase kalibrasi dilakukan untuk memastikan validitas internal. Kegiatan meliputi: (a)~deployment InvenioRDM pada kedua lingkungan; (b)~verifikasi baseline bersih; (c)~pengukuran awal konsistensi; (d)~verifikasi penggunaan sumber daya node di bawah 70\%; dan (e)~uji coba awal untuk memvalidasi protokol pengukuran dan skrip otomasi.
    
    \item \textbf{Eksperimen dan Pengumpulan Data (Oktober -- November 2026)}
    
    Pelaksanaan tiga eksperimen sesuai desain pada Bab III. E1: konsistensi deployment (n = 5 per skenario). E2: \textit{detection time}, \textit{recovery time}, dan \textit{root cause identification time} (n = 10 per skenario). E3: \textit{traceability} operasional (deskriptif, dokumentasi kelengkapan audit log per skenario). Data dikumpulkan dalam bentuk \textit{timestamp}, status aplikasi, dan log. Skenario diacak untuk setiap lingkungan untuk memitigasi \textit{order effect}. Seluruh prosedur dijalankan menggunakan skrip otomasi untuk konsistensi.
    
    \item \textbf{Analisis Statistik dan Interpretasi (November -- Desember 2026)}
    
    Analisis data meliputi: statistik deskriptif (mean, median, rentang, deviasi standar), uji perbandingan (\textit{Mann--Whitney U test}), perhitungan ukuran efek (Cohen's \textit{d}), dan interpretasi hasil. Pembahasan mencakup identifikasi skenario drift yang paling berdampak, faktor yang mempengaruhi hasil, serta ancaman terhadap validitas internal dan eksternal penelitian.
    
    \item \textbf{Penulisan Laporan (Juni 2026 -- Januari 2027)}
    
    Penulisan laporan dilakukan secara paralel sepanjang penelitian. Setiap tahap menghasilkan bab atau bagian yang sesuai. Draft proposal diselesaikan pada pertengahan semester, draft skripsi lengkap pada akhir penelitian.
    
    \item \textbf{Seminar / Sidang (Januari 2027)}
    
    Presentasi dan sidang tugas akhir di depan dosen penguji, mencakup paparan hasil eksperimen, analisis statistik, dan kesimpulan.
\end{enumerate}